\ulohahead{specializace UI-SU}{Prohledávání: neinformované a informované (A*)}
\begin{zadani}
\begin{enumerate}[label=(\alph*)]
  \item \bod{1} Popište formulaci úlohy prohledáváním (stav, operátory, cíl, cena). Rozdíl mezi stromovým a grafovým prohledáváním. Uveďte neinformované metody BFS, DFS a uniform-cost.
  \item \bod{2} Popište algoritmus A* a funkci $f(n)=g(n)+h(n)$. Co je přípustná a co konzistentní heuristika?
  \item \bod{3} Za jakých podmínek je A* optimální? Co znamená dominance heuristik a proč je lepší ,,větší'' přípustná heuristika?
\end{enumerate}
\end{zadani}

\begin{reseni}
\textbf{(a)} \emph{Formulace}: počáteční stav, množina \emph{operátoru} (akce s cenou) generujících následníky, \emph{cílový test}, cena cesty. \emph{Stromové} prohledávání neeviduje navštívené stavy (může opakovat), \emph{grafové} drží uzavřenou množinu (closed set). \emph{BFS}: po vrstvách (fronta), úplné, optimální při jednotkových cenách. \emph{DFS}: do hloubky (zásobník), paměťově levné, neúplné/nesuboptimální. \emph{Uniform-cost}: rozšiřuje uzel s nejmenší $g$ (Dijkstra), optimální pro nezáporné ceny.

\textbf{(b)} \emph{A*} rozšiřuje uzel s nejmenším $f(n)=g(n)+h(n)$, kde $g$ je cena z počátku a $h$ \emph{heuristický} odhad ceny do cíle. Heuristika je \emph{přípustná}, pokud nikdy nepřeceňuje ($0\le h(n)\le h^*(n)$); \emph{konzistentní} (monotónní), pokud $h(n)\le c(n,n')+h(n')$ pro každý přechod (trojúhelníková nerovnost).

\textbf{(c)} A* je optimální se \emph{přípustnou} heuristikou ve stromovém prohledávání a s \emph{konzistentní} heuristikou v grafovém (jinak je nutné znovuotevírání uzlu). \emph{Dominance}: $h_2$ dominuje $h_1$, je-li $h_2(n)\ge h_1(n)$ pro všechny $n$ (a obě přípustné); pak A* s $h_2$ rozšíří méně uzlu (efektivnější), protože vyšší přípustný odhad ořeže víc.
\end{reseni}

\kalibrace{Prohledávání (~15\,\%) je opakovaná ,,Základy UI'' otázka (A* se objevil). Definice A* + přípustnost/konzistence = 2-bodová záloha pro AI slot specializace.}
\past{Přípustná $\neq$ konzistentní (konzistence je silnější a implikuje přípustnost). V grafovém A* je pro optimalitu potřeba konzistence, ne jen přípustnost.}
